\section{Completeness: from PROs to PROPs}
\label{sec:completeness}

For each fragment~\(\bullet\) of \cref{sec:circuits-relations} we have a presented PROP
\(\QCirc_\bullet=\cat{P}_\bullet/\QC_\bullet\) and a strict symmetric monoidal interpretation
\(\interp{\cdot}_\bullet:\QCirc_\bullet\to\cat{Qubit}_\bullet\) (or
\(\cat{Qutrit}_{\mathit{Cliff3}}\) in the qutrit case).  Completeness means faithfulness of this
interpretation.  Complete equational theories for these fragments already exist in the
literature
\cite{SelingerStabilizer,SelingerRealStabilizer,qutritclifford,SelingerCliffordPlusT,SelingerCliffordPlusCS,AmyCNOTDihedral}.
Those theorems supply the source completeness results: their normal forms prove faithfulness for
older source presentations, and the transfer argument in this section checks that the source syntax
and our smaller figure-defined target syntax present the same semantic subPROP.

The transfer has to align three conventions before the cited normal forms can be used for our
quotient PROPs.  If a source theorem is stated for a PRO, \cref{lem:fromprotoprop} adds the
structural symmetries needed to read it as a PROP theorem.  For qutrit Clifford and
Clifford\(+CS\), the source and target scalar subgroups differ: the source presentations make fewer
global phases explicit, which is harmless for the original proofs but does not match the standard
strict matrices fixed in \cref{sec:circuits-relations}.  The scalar-refinement step adjoins the
missing roots of unity, so source and target have the same strict semantics.  Encoding and decoding
are then identity-on-objects PROP morphisms, in the sense of \cref{def:prop-morphism}, that compare
the source generators and relations with ours; once these maps preserve semantics and decode source
axioms to derivable target equations, \cref{th:completeness} transports faithfulness.

\subsection{Scalar refinement}
\label{subsec:scalar-refinement}

Because our semantics uses strict unitary equality rather than projective equality, global phases
are visible.  Equivalently, the scalar subgroup \(S(\cat{C})=\cat{C}(0,0)\) of the semantic PROP is
part of the data that must match between source and target presentations.
For our purposes, a scalar mismatch is therefore a mismatch between presentation conventions: the
source convention makes fewer global phases explicit, while our convention keeps the standard gate
matrices and records the extra phases as scalars.

Only two source completeness results require scalar adjustment.  The qutrit Clifford presentation of
\cite{qutritclifford} uses an order-\(6\) scalar subgroup whereas our semantics uses order \(12\), and
the Clifford\(+CS\) presentation of \cite{SelingerCliffordPlusCS} uses order \(4\) whereas our qubit
semantics uses order \(8\).  In both cases we appeal to the scalar-refinement construction of
Appendix~\ref{app:scalar-refinement}, which adjoins the missing scalar generator and is conservative
by \cref{lem:scalar-refinement}.  Source and target can then be compared inside the same semantic
subPROP, while the refinement remains conservative on the old source signature.

\subsection{Encoding, decoding, and transfer of completeness}

The five fragments handled by encoding and decoding reduce to one generator-and-relation
comparison.

Fix one encoding/decoding transfer instance.  Here \(\cat{P}\) and \(\mathcal{R}\) denote the
target free PROP and target relations, while \(\cat{P}^{\mathrm{src}}\) and
\(\mathcal{R}^{\mathrm{src}}\) denote the corresponding source free PROP and source relations,
after PRO-to-PROP lifting and scalar refinement where those steps apply.

\begin{definition}\label{def:encdec}
An encoding/decoding pair consists of a target-to-source PROP morphism
\(E:\cat{P}\to\cat{P}^{\mathrm{src}}\) and a source-to-target PROP morphism
\(D:\cat{P}^{\mathrm{src}}\to\cat{P}\) such that
\(\interp{C}=\interp{E(C)}_{\mathrm{src}}\) for every target circuit \(C\in\cat{P}\) and
\(\interp{C'}_{\mathrm{src}}=\interp{D(C')}\) for every source circuit
\(C'\in\cat{P}^{\mathrm{src}}\).
\end{definition}

After importing a source completeness theorem, the comparison is local: it involves only the chosen
generators and source relations.  A typical mismatch is that a source presentation takes
controlled-\(Z\) as primitive, whereas our target signature takes CNOT.  The decoder replaces source
controlled-\(Z\) by its fixed expansion through CNOT and Hadamards; the target-generator check then
proves, for example, that decoding the encoded CNOT derives CNOT again, while the source-relation
check verifies each source equation after the same replacement.  The non-identity parts of the maps
used here are recorded in \cref{tab:ED-clifford-family}.

The transfer therefore has a finite checklist: the round-trip equality \(D(E(g))=g\) for each
target generator \(g\), and the decoded equality \(D(L)=D(R)\) for each source axiom \(L=R\).
This is the presentation-level analogue of Reidemeister--Schreier rewriting
\cite{Reidemeister1927,Schreier1927}: instead of computing a subgroup presentation, we compare two
chosen presentations by following the images of their generators and defining relations.

\begin{lemma}\label{lem:decenc}
Assume \(\cat{P}\) is presented by generators \(\Sigma\) and relations \(\mathcal{R}\).
If for every non-structural generator \(g\in\Sigma\),
\(
  \cat{P}/\mathcal{R}\vdash D(E(g)) = g,
\)
then for every circuit \(C\) built from \(\Sigma\),
\(
  \cat{P}/\mathcal{R}\vdash D(E(C)) = C.
\)
\end{lemma}
\begin{proof}
By structural induction on \(C\), using that \(E\) and \(D\) preserve \(\circ\), \(\otimes\), and the
structural morphisms.
\end{proof}

\begin{lemma}\label{lem:decrelations}
Assume that for every axiom \(L=R\) in \(\mathcal{R}^{\mathrm{src}}\),
\(
  \cat{P}/\mathcal{R}\vdash D(L) = D(R).
\)
Then \(D\) induces a PROP morphism
\(
  \overline{D}:\cat{P}^{\mathrm{src}}/\mathcal{R}^{\mathrm{src}}\to\cat{P}/\mathcal{R},
\)
and any derivable equality in the source quotient transports to a derivable equality between its
decodings in the target theory.
\end{lemma}
\begin{proof}
The hypothesis says exactly that \(D\) equalises all source relations, hence factors through the
quotient.
\end{proof}

\begin{theorem}\label{th:completeness}
Assume \(\cat{P}^{\mathrm{src}}/\mathcal{R}^{\mathrm{src}}\) is complete for
\(\interp{\cdot}_{\mathrm{src}}\), and let \((E,D)\) be an encoding/decoding pair satisfying the
hypotheses of \cref{lem:decenc,lem:decrelations}.  Then \(\cat{P}/\mathcal{R}\) is complete for
\(\interp{\cdot}\).
\end{theorem}
\begin{proof}
If target circuits \(C_1,C_2\) have the same semantics, then semantic preservation of \(E\) gives
\(\interp{E(C_1)}_{\mathrm{src}}=\interp{E(C_2)}_{\mathrm{src}}\).  Source completeness yields
\(E(C_1)=E(C_2)\) in
\(\cat{P}^{\mathrm{src}}/\mathcal{R}^{\mathrm{src}}\); applying the quotient morphism
\(\overline{D}\) from \cref{lem:decrelations} gives \(D(E(C_1))=D(E(C_2))\) in
\(\cat{P}/\mathcal{R}\), and \cref{lem:decenc} rewrites both sides to \(C_1=C_2\).
\end{proof}

The same proof gives the bounded form used in \cref{cor:clifford-family-complete}: if the source
interpretation is faithful on all hom-sets \(n\to m\) with \(n,m\le k\), then the target
interpretation is faithful on the same range, because \(E\) and \(D\) are identity-on-objects PROP
morphisms.

\subsection{Fragment-wise completeness}

Applying the transfer to the five encoding/decoding fragments and the direct CNOT-dihedral
argument gives the completeness bounds used by the minimality theorem.

\begin{theorem}\label{cor:clifford-family-complete}
\begin{enumerate}
\item \(\QCirc_{\mathit{Cliff}}\) is complete for \(\cat{Qubit}_{\mathit{Cliff}}\).
\item \(\QCirc_{\mathit{RCliff}}\) is complete for \(\cat{Qubit}_{\mathit{RCliff}}\).
\item \(\QCirc_{\mathit{Cliff3}}\) is complete for \(\cat{Qutrit}_{\mathit{Cliff3}}\).
\item \(\QCirc_{\mathit{CliffT}}\) is complete up to \(2\) qubits for \(\cat{Qubit}_{\mathit{CliffT}}\).
\item \(\QCirc_{\mathit{CliffCS}}\) is complete up to \(3\) qubits for \(\cat{Qubit}_{\mathit{CliffCS}}\).
\item \(\QCirc_{\mathit{CNOTdihe}}\) is complete for \(\cat{Qubit}_{\mathit{CNOTdihe}}\).
\end{enumerate}
\end{theorem}

\begin{proof}
For the five fragments using encoding/decoding, start from the corresponding cited source
completeness theorem; in the qutrit Clifford and Clifford\(+CS\) cases, first pass to the
scalar-refined source presentation supplied by \cref{lem:scalar-refinement}.  The arity bounds in
the Clifford+T and Clifford+CS items are the bounds inherited from those source results.

The appendices record the data needed to invoke \cref{th:completeness}: the earlier rule sets in
\cref{appendix:oldgraphicallanguage1,appendix:oldgraphicallanguage2,appendix:oldgraphicallanguage3,appendix:oldgraphicallanguage4,appendix:oldgraphicallanguage5,appendix:oldgraphicallanguage6},
the encoding/decoding maps in \cref{tab:ED-clifford-family}, the generator checks for
\cref{lem:decenc} in Appendix~\ref{app:proofs-lem1}, and the source-axiom checks for
\cref{lem:decrelations} in Appendix~\ref{appendix:decodage}.  With these ingredients fixed, the
remaining derivability obligations are exactly the finite checks of the transfer lemmas; no further
normal-form argument is used here.

For CNOT-dihedral, use the completeness theorem of \cite{AmyCNOTDihedral} directly: no
encoding/decoding is needed, and Appendix~\ref{app:CNOT-dihedral} derives, in \(\QCdihedral\),
each source axiom removed from our simplified presentation.
\end{proof}
